

\subsection{冻结参数，把预算分给多次尝试}
\phantomsection\label{concept:test-time-compute}

\textbf{把计算预算分给更多尝试，再用验证器筛选。}
前面的奖励进入优化器，改变未来请求使用的参数；现在冻结参数，
把同一个验证器用于本次候选的筛选。采样与验证只改变运行状态。

沿用 RLVR 的 $17+25$ 算例，验证器直接执行加法，正确答案是 $42$。
现在\textbf{冻结参数}，让模型生成 $n$ 份完整候选，再对每份运行验证器，
选出通过验证的答案；全部未通过就报告本次未解出。
这就是用完美二值验证器做 Best-of-$n$ 的最小版本。
一般的 Best-of-$n$ 也可按奖励模型的分数选最高项，其准确性取决于评分器。

\begin{example}[breakable=false]
\noindent\emph{从一次尝试到八次：先算全部失败的概率}\par
为教学假设每次\textbf{独立}采样都从 $\{41,42,43,50\}$ 中等概率取一个答案，
答对的概率固定为 $p=0.25$。
一次失败概率是 $0.75$；$n$ 次全失败概率是 $0.75^n$。
因此至少有一份正确候选的概率为
\[
P_n=1-(1-p)^n.
\]
\[
\begin{array}{c|ccccc}
n&1&2&4&8&16\\\hline
P_n&0.2500&0.4375&0.6836&0.8999&0.9900
\end{array}
\]
完美验证器能选出存在的正确答案，所以这里的最终成功率也等于 $P_n$。
从 $4$ 次增到 $8$ 次提高约 $21.63$ 个百分点，
从 $8$ 次增到 $16$ 次只再提高约 $9.01$ 个百分点：预算继续翻倍，收益逐渐变小。
\end{example}

\par\medskip
\Needspace{85mm}
\begin{center}
\tagstructbegin{tag=Figure,alttext={Best-of-N 教学成功率图。横轴为一到十六份候选，纵轴为交付正确答案的概率。独立候选且完美核验时，单次成功率四分之一，曲线经过 N 为一、二、四、八、十六时的百分之二十五、四十三点七五、六十八点三六、八十九点九九、九十九点零零。候选完全复制时，成功率始终为百分之二十五。}}
\begin{tikzpicture}
\begin{axis}[width=.88\linewidth,height=61mm,
 xmin=1,xmax=16,ymin=0,ymax=1.06,
 xtick={1,2,4,8,16},ytick={0,.25,.5,.75,1},
 xlabel={候选数 $N$},ylabel={成功率},
 tick label style={font=\footnotesize},label style={font=\small},
 axis lines=left,axis line style={draw=muted},
 grid=major,grid style={draw=linegray},
 legend style={font=\footnotesize,draw=none,fill=white,at={(.97,.18)},anchor=south east}]
\addplot[blue,very thick,domain=1:16,samples=100]{1-.75^x};
\addlegendentry{独立采样＋完美核验}
\addplot[only marks,mark=*,blue] coordinates {(1,.25) (2,.4375) (4,.68359375) (8,.89988708496) (16,.98997740424)};
\addlegendentry{手算的五个点}
\addplot[orange,thick,dashed,domain=1:16]{.25};
\addlegendentry{完全复制同一次答案}
\end{axis}
\end{tikzpicture}
\captionof{figure}{多尝试的收益取决于候选的差异和核验能力；曲线是教学假设下的计算结果}
\label{fig:best-of-n-success}
\tagstructend
\end{center}


\textbf{现在动手：至少需要多少次，才能达到 $90\%$？}
由 $1-0.75^n\geq0.9$，得到
\[
n\geq\frac{\ln0.1}{\ln0.75}\approx8.004,
\qquad n_{\min}=9.
\]
八次的 $89.9887\%$ 虽可四舍五入成 $90\%$，仍低于题设门槛；
九次为约 $92.49\%$。
若每份候选固定生成 $128$ 个 Token，完整采样八份就是 $1024$ 个输出 Token，
还要加上验证与输入处理开销；并行执行可缩短等待，总运算与资源占用照样计入。

\begin{boundarybox}
$1-(1-p)^n$ 描述\textbf{固定单次成功率、独立候选}中至少一次答对的概率。
若八份候选完全复制同一次随机答案，成功率仍只有 $25\%$。
若验证器误把与 $42$ 相差 $1$ 的 $41$、$43$ 也判为正确，并交付第一份通过的候选：
至少一份通过的概率是 $1-0.25^n$，而第一份通过者在 $41,42,43$ 中等可能，
\[
P^{\rm faulty}_n=\frac{1-0.25^n}{3}:\quad
P^{\rm faulty}_1=0.25,\ P^{\rm faulty}_2=0.3125,\ P^{\rm faulty}_4\approx0.3320,\ P^{\rm faulty}_8\approx0.3333.
\]
再多采样也停在约 $1/3$：“候选中存在正确解”与“最后交付正确解”要分开评测。
在开放任务中，模型自评分、程序测试与资料出处分别有自己的覆盖范围，
验证器需要独立检查。
\end{boundarybox}

RLVR 把验证奖励用于\textbf{更新参数}；这里的 Best-of-$n$ 用验证结果\textbf{筛选本次候选}。
串行修正和搜索也能使用预算，实际策略要在相同任务和预算下比较；
可参阅 \href{https://arxiv.org/abs/2408.03314}{Snell 等对测试时计算的研究}。
\textbf{复算入口：}\texttt{python3 scripts/test\_time\_compute\_demo.py}，
脚本枚举小样本空间，分别核对独立采样、完全相关采样与错误验证器的结果。


\subsection{“让模型写步骤”与“训练模型推理”不是一回事}
\phantomsection\label{concept:cot-prompt}

\textbf{写出中间步骤，为后续 Token 提供一个可读取的工作区。}
思维链提示（Chain-of-Thought, CoT）让模型先写计算过程，再给结果。
每个新 Token 都经过一次模型前向计算，后续生成又能读取此前写出的步骤，
于是多步任务得到更多串行计算与中间记录。

CoT 提示只改变输入，训练才修改参数，推理时搜索则额外花费计算；
三者的分工集中列在本章末的\hyperref[sec:training-inference-roles]{分工表}中。

\begin{boundarybox}
步骤写得完整仍需逐项核对；它们是生成文本，无法直接证明模型内部采用了同一推理过程。
错误步骤会向后传播，额外生成也会增加 Token 与延迟成本。
训练后的收益用独立任务测试和失败案例衡量。
\end{boundarybox}

\paragraph{回到季度报告任务}
Few-shot 给格式示例，SFT 用核对过的示范训练格式与缺证据时的回应，偏好反馈比较来源与单位是否清楚，
可验证奖励检查数值、公式与出处。读取报告的权限与原表证据仍由应用提供；
\hyperref[intro:23d-agent-systems]{检索与任务系统}一章将接着说明证据怎样进入上下文。


\subsection{不同训练与推理策略的分工}
\phantomsection\label{sec:training-inference-roles}
提示、训练与推理时计算的区别全书只在这里列一次：
提示策略改变输入而不修改参数；SFT、偏好与奖励信号修改参数；
推理时采样、验证或搜索在固定参数下多花计算，换取一定条件下更高的成功率。

\begin{center}
\small
\renewcommand{\arraystretch}{1.3}
\tagpdfsetup{table/header-rows={1}}
\begin{tabularx}{\linewidth}{@{}p{30mm}X@{}}
\toprule
阶段/方法 & 主要改变什么\\
\midrule
预训练 & 大规模数据中学习一般 Token 预测模式\\
SFT & 利用示范答案塑造特定任务与回应方式\\
RLHF/DPO 等偏好优化 & 根据比较或反馈优化回答偏好；通常约束新策略不远离参考策略\\
RLVR & 利用可程序化检查的奖励训练某些任务能力\\
CoT 提示 & 改变输入组织方式，不必修改模型权重\\
推理时搜索/验证 & 使用额外计算与外部规则筛选候选结果\\
\bottomrule
\end{tabularx}
\end{center}

\chapterchecktasks{23c-post-training}{理解、手算与判断}{
\label{check:23cc-test-time-compute}
\item\textbf{比较：}SFT、成对偏好学习与可验证奖励的反馈信号有何不同？各举一个适合的任务。\chapterinlinehint{先列出三种方法各自实际收到的反馈。}
\item\textbf{手算：}SFT 目标 Token 的概率为 $0.25$，求单位置交叉熵；若四个组内奖励是 $(1,0,1,0)$，按总体标准差标准化得到哪些相对优势？\chapterinlinehint{先算负对数，再用奖励减均值并除标准差。}
\item\textbf{{\symboliconfont ★} 找错：}同学把这四个相对优势直接当成完整 GRPO 更新，遗漏了什么策略目标与数值边界？\chapterinlinehint{在相对优势之后继续写策略目标，并检查零方差。}
\item\textbf{解释：}向同学解释 RLVR 与 Best-of-$n$ 各把验证结果用在何处。\chapterinlinehint{检查本次是否更新模型权重。}
\item\textbf{估算：}若每份候选独立答对概率为 $1/2$，完美验证器检查四份，最终成功率是多少？\chapterinlinehint{先求所有候选都失败的概率。}
\item\textbf{复现：}写十行以内 Python，比较四份独立候选与四份完全相同候选的成功率。解释两者为什么不同。\chapterinlinehint{分别枚举独立样本空间与共享随机变量。}
}
\par\smallskip
训练改变参数；测试时计算把固定模型的预算分给生成、搜索和验证。
先报告单次成功率，再报告候选相关性、验证质量与总成本，才能判断多算是否值得。
